\expandafter\ifx\csname CAPETINCLUDEMODE\endcsname\relax
\documentclass[11pt,a4paper]{article}
% preamble.tex — packages et commandes communes a toutes les fiches CAPET SII IE
\input{glyphtounicode}
\pdfgentounicode=1
\usepackage{cmap}
\usepackage[T1]{fontenc}
\usepackage[utf8]{inputenc}
\usepackage[french]{babel}
\usepackage[a4paper, top=1.6cm, bottom=1.6cm, left=1.4cm, right=1.4cm]{geometry}
\setlength{\headheight}{14.5pt}
\usepackage{amsmath, amssymb, amsfonts}
\usepackage{siunitx}
\sisetup{locale=FR, output-decimal-marker={,}, inter-unit-product=\ensuremath{{}\cdot{}}}
\usepackage{array, booktabs, multirow, tabularx}
\usepackage{xcolor}
\usepackage[most, breakable]{tcolorbox}
\usepackage{multicol}
\usepackage{enumitem}
\usepackage{mathtools}
\usepackage{lmodern}
\usepackage{microtype}
\usepackage{fancyhdr}
\usepackage{tikz}
\usetikzlibrary{positioning,arrows.meta}
\usepackage{lastpage}

% ── Couleurs par domaine ──────────────────────────────────────────────
\definecolor{colElec}{RGB}{26, 82, 148}       % Électrotechnique — bleu foncé
\definecolor{colElecLight}{RGB}{210, 225, 245} % fond bleu clair
\definecolor{colAuto}{RGB}{34, 120, 66}        % Automatique — vert
\definecolor{colMeca}{RGB}{185, 90, 20}        % Mécanique — orange
\definecolor{colTherm}{RGB}{160, 40, 40}       % Thermique — rouge
\definecolor{colQT}{RGB}{232, 244, 233}        % fond question type (vert très léger)
\definecolor{colJury}{RGB}{255, 243, 205}      % fond jury (jaune ambre)
\definecolor{colJuryFrame}{RGB}{200, 130, 0}   % bordure jury
\definecolor{colGrayLight}{RGB}{245, 245, 247} % fond gris neutre
\definecolor{colSeparator}{RGB}{180, 180, 180} % gris séparateur

% ── Style tcolorbox de section principale ────────────────────────────
\tcbset{
  ficheSection/.style={
    breakable,
    enhanced,
    colback=#1!7!white,
    colframe=#1,
    colbacktitle=#1,
    coltitle=white,
    fonttitle=\bfseries\normalsize,
    left=4pt, right=4pt, top=3pt, bottom=4pt,
    before skip=6pt, after skip=6pt,
  }
}

% ── Style encadré question type ──────────────────────────────────────
\tcbset{
  questiontype/.style={
    breakable,
    enhanced,
    colback=colQT,
    colframe=colElec!55!black,
    colbacktitle=colElec!55!black,
    coltitle=white,
    fonttitle=\bfseries\small,
    left=4pt, right=4pt, top=3pt, bottom=3pt,
    before skip=5pt, after skip=5pt,
  }
}

% ── Style encadré jury ───────────────────────────────────────────────
\tcbset{
  jury/.style={
    breakable,
    enhanced,
    colback=colJury,
    colframe=colJuryFrame,
    fonttitle=\bfseries\normalsize,
    left=4pt, right=4pt, top=3pt, bottom=3pt,
  }
}

% ── Mise en page ─────────────────────────────────────────────────────
\pagestyle{fancy}
\fancyhf{}
\renewcommand{\headrulewidth}{0.4pt}
\renewcommand{\footrulewidth}{0.3pt}

\setlist[itemize]{leftmargin=1.2em, itemsep=1pt, topsep=2pt, parsep=0pt}
\setlist[enumerate]{leftmargin=1.2em, itemsep=3pt, topsep=2pt, parsep=0pt}

\setlength{\parskip}{2pt}
\setlength{\columnsep}{10pt}

\def\CAPETBEGINFILE{\begin{document}}
\def\CAPETENDFILE{\end{document}}
\else
\def\CAPETBEGINFILE{}
\def\CAPETENDFILE{}
\fi

\fancyhead[L]{\textbf{\textcolor{colAuto}{T28}} - Logique combinatoire}
\fancyhead[R]{\small\textit{CAPET SII IE - Session 2026}}
\fancyfoot[C]{\small\thepage\,/\,\pageref{LastPage}}

\CAPETBEGINFILE

% -----------------------------------------------------------------------
%  BLOC EN-TeTE
% -----------------------------------------------------------------------

\begin{tcolorbox}[
  enhanced, breakable,
  colback=colAuto,
  colframe=colAuto,
  left=6pt, right=6pt, top=5pt, bottom=5pt,
  before skip=0pt, after skip=4pt
]
  {\color{white}\bfseries\large T28 \quad Logique combinatoire}\\[2pt]
  {\color{colAuto!30!white}\small
    \textbf{Domaine :} electronique numerique - Logique et codage \quad
    \textbf{Presence :} \textbf{4/13} sujets (implicite) - theme de \emph{culture}
  }\\[1pt]
  {\color{colAuto!20!white}\footnotesize
    Sous-themes : algebre de Boole - portes logiques - De Morgan - table de verite - Karnaugh - circuits combinatoires - codage binaire / BCD / Gray
  }
\end{tcolorbox}

% -----------------------------------------------------------------------
%  1. RAPPELS DE COURS
% -----------------------------------------------------------------------

\begin{tcolorbox}[
  ficheSection=colAuto,
  title={1. Rappels de cours},
]

% -- 1.1 Operateurs fondamentaux --------------------------------------
\textbf{1.1 Algebre de Boole - operateurs et theoremes}

\smallskip
\begin{center}
\renewcommand{\arraystretch}{1.25}
\begin{tabular}{llll}
\toprule
\textbf{Operation} & \textbf{Notation} & \textbf{Identites notables} & \\
\midrule
ET (AND) & $A \cdot B$ & Idempotence : $A \cdot A = A$ & \\
OU (OR) & $A + B$ & Absorption : $A(A+B)=A$ & \\
NON (NOT) & $\overline{A}$ & Complementarite : $A \cdot \overline{A} = 0$ & \\
\midrule
\multicolumn{4}{l}{\textbf{Theoremes de De Morgan :}} \\
\multicolumn{4}{l}{$\overline{A \cdot B} = \overline{A} + \overline{B}$ \hspace{2em} $\overline{A + B} = \overline{A} \cdot \overline{B}$} \\
\bottomrule
\end{tabular}
\end{center}

\medskip\hrule\vspace{4pt}

% -- 1.2 Portes logiques -----------------------------------------------
\textbf{1.2 Portes logiques elementaires (2 entrees)}

\begin{center}
\renewcommand{\arraystretch}{1.2}
\begin{tabular}{lll}
\toprule
\textbf{Porte} & \textbf{Expression} & \textbf{Remarque} \\
\midrule
AND & $S = A \cdot B$ & $S = 1$ ssi $A = B = 1$ \\
OR & $S = A + B$ & $S = 0$ ssi $A = B = 0$ \\
NOT & $S = \overline{A}$ & Inverseur \\
NAND & $S = \overline{A \cdot B}$ & Porte universelle \\
NOR & $S = \overline{A + B}$ & Porte universelle \\
XOR & $S = A \oplus B$ & $S = 1$ ssi $A \neq B$ \\
XNOR & $S = \overline{A \oplus B}$ & $S = 1$ ssi $A = B$ \\
\bottomrule
\end{tabular}
\end{center}

\textbf{Realisation universelle :} toute fonction est realisable avec des NAND seuls (ou NOR seuls).

\medskip\hrule\vspace{4pt}

% -- 1.3 Table de verite et formes canoniques --------------------------
\textbf{1.3 Table de verite - formes canoniques}

Pour $n$ variables : $2^n$ lignes.
\begin{itemize}
  \item \textbf{SOP (Somme de mintermes)} : somme des mintermes pour lesquels $S = 1$
  \item \textbf{POS (Produit de maxtermes)} : produit des maxtermes pour lesquels $S = 0$
\end{itemize}

\textit{Exemple - demi-additionneur :}
\begin{center}
\renewcommand{\arraystretch}{1.1}
\begin{tabular}{cccc}
\toprule
$A$ & $B$ & $S = A \oplus B$ & $C = A \cdot B$ \\
\midrule
0 & 0 & 0 & 0 \\
0 & 1 & 1 & 0 \\
1 & 0 & 1 & 0 \\
1 & 1 & 0 & 1 \\
\bottomrule
\end{tabular}
\end{center}

\medskip\hrule\vspace{4pt}

% -- 1.4 Tableau de Karnaugh -------------------------------------------
\textbf{1.4 Simplification par tableau de Karnaugh}

\textbf{Regles :} grouper les - 1 - adjacents par puissances de 2 (groupes de 2 - $-1$ var, de 4 - $-2$ var, de 8 - $-3$ var). Le tableau est \textbf{torique} : les bords sont adjacents.

\medskip\hrule\vspace{4pt}

% -- 1.5 Circuits combinatoires courants ------------------------------
\textbf{1.5 Circuits combinatoires courants}

\begin{center}
\renewcommand{\arraystretch}{1.2}
\begin{tabular}{lp{9cm}}
\toprule
\textbf{Circuit} & \textbf{Fonction} \\
\midrule
Multiplexeur $2^n \to 1$ & Selectionne 1 entree parmi $2^n$ selon $n$ bits de selection \\
Decodeur $n \to 2^n$ & Active une sortie parmi $2^n$ (adressage memoire) \\
Additionneur complet & $S = A \oplus B \oplus C_{in}$ ; $C_{out} = A\cdot B + C_{in}(A\oplus B)$ \\
Comparateur & Genere $A > B$, $A = B$, $A < B$ \\
\bottomrule
\end{tabular}
\end{center}

\medskip\hrule\vspace{4pt}

% -- 1.6 Codage numerique ----------------------------------------------
\textbf{1.6 Codages numeriques}

\begin{itemize}
  \item \textbf{Binaire naturel :} $N_{10} = b_{n-1}\cdot 2^{n-1} + \cdots + b_1 \cdot 2 + b_0$
  \item \textbf{BCD :} chaque decimale codee sur 4 bits (ex. : $9 = 1001_2$)
  \item \textbf{Code Gray :} un seul bit change entre deux mots consecutifs - codeurs angulaires
  \item \textbf{Hexadecimal :} base 16 ; chiffres 0-9 puis AeF ; $0\text{x}45 = 69_{10}$
\end{itemize}

\textit{Connexion sujets CAPET (2023) :} codeur incremental 12 bits, resolution $0{,}15$ mm/impulsion ; 8 bits de poids fort retenus (resolution effective $0{,}15 \times 2^4 = 2{,}4$ mm).

\end{tcolorbox}

% -----------------------------------------------------------------------
%  2. FORMULAIRE SYNTHeTIQUE
% -----------------------------------------------------------------------

\begin{tcolorbox}[
  ficheSection=colAuto,
  title={2. Formulaire synthetique},
]
\small
\begin{multicols}{2}
\renewcommand{\arraystretch}{1.45}
\begin{tabular}{@{}p{5cm}l@{}}
\toprule
\textbf{Identite / formule} & \textbf{Nom} \\
\midrule
$\overline{A \cdot B} = \overline{A} + \overline{B}$ & De Morgan AND \\
$\overline{A + B} = \overline{A} \cdot \overline{B}$ & De Morgan OR \\
$A(A+B) = A$ & Absorption \\
$A+A\cdot B = A$ & Absorption \\
$A \oplus B = A\overline{B} + \overline{A}B$ & XOR developpe \\
\bottomrule
\end{tabular}

\columnbreak

\renewcommand{\arraystretch}{1.45}
\begin{tabular}{@{}p{5cm}l@{}}
\toprule
\textbf{Formule} & \textbf{Usage} \\
\midrule
$S = A \oplus B \oplus C_{in}$ & Additionneur (somme) \\
$C_{out} = AB + C_{in}(A\oplus B)$ & Additionneur (retenue) \\
$n$ variables - $2^n$ lignes & Table verite \\
$q = U_{PE}/2^n$ & Quantum CAN \\
$f_e \geq 2\,f_{max}$ & Shannon \\
\bottomrule
\end{tabular}
\end{multicols}
\normalsize
\end{tcolorbox}

% -----------------------------------------------------------------------
%  3. QUESTIONS TYPES
% -----------------------------------------------------------------------

\begin{tcolorbox}[
  ficheSection=colAuto,
  title={3. Questions types - extraites des sujets reels},
]

% -- QT1 -------------------------------------------------------------

% Mini-procedure de resolution type copie concours
\begin{tcolorbox}[questiontype,title={Procedure type copie concours (5 etapes)}]
\begin{enumerate}[leftmargin=1.5em]
  \item \textbf{Qualifier le modele} : hypotheses, domaine de validite, grandeurs d'entree/sortie.
  \item \textbf{Ecrire les lois} : equations de base avec symboles et unites SI.
  \item \textbf{Isoler l'inconnue} : expression litterale, puis application numerique tracee.
  \item \textbf{Controler l'ordre de grandeur} : coherence physique et unite finale.
  \item \textbf{Conclure concours} : phrase courte orientee decision technique.
\end{enumerate}
\end{tcolorbox}

% Visuel de methode : chemin de resolution CAPET
\begin{center}
\begin{tikzpicture}[font=\small, node distance=7mm and 6mm, >=Stealth]
  \node[draw, rounded corners, fill=colElec!8, inner sep=4pt] (n1) {Notions};
  \node[draw, rounded corners, fill=colElec!8, right=of n1, inner sep=4pt] (n2) {Methode};
  \node[draw, rounded corners, fill=colElec!8, right=of n2, inner sep=4pt] (n3) {Application annale};
  \node[draw, rounded corners, fill=colElec!8, right=of n3, inner sep=4pt] (n4) {Verification jury};
  \draw[->, thick, colElec] (n1) -- (n2);
  \draw[->, thick, colElec] (n2) -- (n3);
  \draw[->, thick, colElec] (n3) -- (n4);
\end{tikzpicture}
\end{center}
\begin{tcolorbox}[questiontype,
  title={QT1 - Codage binaire / hexadecimal \hfill\normalfont\footnotesize E2 disciplinaire 2023 (D2-D3)}
]
\textbf{enonce.}
Un codeur incremental genere une impulsion tous les $0{,}15$ mm. La distance maximale - coder est $680$ mm. Quel est le nombre minimal de bits necessaires - Le registre ne retient que les 8 bits de poids fort. Quelle est la resolution effective de position - 
\tcbline

\textbf{Corrige.}

Nombre de positions : $N = 680 / 0{,}15 = 4\,533$. Nombre de bits : $2^{12} = 4\,096 < 4\,533 < 8\,192 = 2^{13}$ - \textbf{13 bits}.

Les 8 bits de poids fort correspondent - ignorer les 5 bits de poids faible :
resolution effective $= 0{,}15 \times 2^5 = 0{,}15 \times 32 = \SI{4{,}8}{\milli\metre}$.
\end{tcolorbox}

\vspace{4pt}

% -- QT2 -------------------------------------------------------------

\begin{tcolorbox}[questiontype,
  title={QT2 - Application de De Morgan \hfill\normalfont\footnotesize Cours programme}
]
\textbf{enonce.}
Simplifier l'expression $S = \overline{\overline{A} \cdot \overline{B}}$ et identifier la porte correspondante. Puis simplifier $F = \overline{A + B + C}$.

\tcbline

\textbf{Corrige.}

$S = \overline{\overline{A} \cdot \overline{B}} = \overline{\overline{A}} + \overline{\overline{B}} = A + B$ \quad - \textbf{porte OR} realisee par un NOR suivi d'un NON.

$F = \overline{A + B + C} = \overline{A} \cdot \overline{B} \cdot \overline{C}$ \quad (De Morgan generalise) - NOR - 3 entrees.
\end{tcolorbox}

\vspace{4pt}

% -- QT3 -------------------------------------------------------------

\begin{tcolorbox}[questiontype,
  title={QT3 - Conversion hexadecimal / decimal \hfill\normalfont\footnotesize E2 disciplinaire 2022 (Q.17)}
]
\textbf{enonce.}
Un registre Siemens S7-1500 contient la valeur $\text{0x}2A$ (hexadecimal). Donner la valeur decimale. Un masquage par ET avec $\text{0x}0F$ est applique. Que conserve-t-on - 
\tcbline

\textbf{Corrige.}

$\text{0x}2A = 2 \times 16 + 10 = 42_{10}$ ; en binaire : $0010\,1010_2$.

Masque $\text{0x}0F = 0000\,1111_2$. Le ET logique donne $0000\,1010_2 = 10_{10}$ (octet de poids faible conserve).
\end{tcolorbox}

\end{tcolorbox}

% -----------------------------------------------------------------------
%  4. ATTENTES DU JURY
% -----------------------------------------------------------------------

\begin{tcolorbox}[
  jury,
  title={4. Attentes du jury - Erreurs recurrentes et points de vigilance},
]

\begin{enumerate}[label=\textbf{\arabic*.}, leftmargin=*, itemsep=5pt]

  \item \textbf{De Morgan - le piege classique.} $\overline{A + B} = \overline{A} \cdot \overline{B}$ (ET), \textbf{pas} $\overline{A} + \overline{B}$. Toujours verifier sur un cas numerique ($A=1$, $B=0$).

  \item \textbf{Tableau de Karnaugh torique.} Les cases des bords gauche et droit sont adjacentes, de meme pour haut et bas. Oublier de boucler est une perte de simplification.

  \item \textbf{Confusion bit / octet.} Le jury 2023 signale : "beaucoup d'amalgames entre bit et octet". 1 octet = 8 bits ; $2^8 = 256$ valeurs.

  \item \textbf{XOR vs OR inclusif.} $A \oplus B = 1$ ssi $A \neq B$ (OR \emph{exclusif}). Le OR standard est $A = 1$ \textbf{ou} $B = 1$ \textbf{ou} les deux.

  \item \textbf{Codage binaire signe.} Pour les registres d'automate sur 16 bits signes (Siemens S7), le MSB est le bit de signe (complement - 2). Plage : $-32\,768$ - $+32\,767$.

  \item \textbf{Contexte CAPET.} La logique combinatoire pure (Karnaugh, SOP/POS) est peu presente dans les epreuves ecrites. Maetriser \emph{surtout} le codage binaire, hexadecimal, et les operations de masquage/decalage qui apparaissent regulierement.

\end{enumerate}

\medskip
\begin{tcolorbox}[
  colback=colJuryFrame!12!colJury,
  colframe=colJuryFrame,
  left=4pt, right=4pt, top=2pt, bottom=2pt,
  fontupper=\small\itshape
]
Rapport jury 2023 : "Les candidats doivent maetriser les bases de la numeration (binaire, hexadecimal, BCD) et savoir effectuer des conversions. De nombreuses erreurs de calcul sur les puissances de 2 sont constatees."
\end{tcolorbox}

\end{tcolorbox}

% -----------------------------------------------------------------------
%  PIED - PReSENCE DANS LES SUJETS
% -----------------------------------------------------------------------

\begin{tcolorbox}[
  enhanced, breakable,
  colback=colGrayLight,
  colframe=colSeparator,
  left=4pt, right=4pt, top=3pt, bottom=3pt,
  before skip=4pt, after skip=0pt,
  fonttitle=\small\bfseries,
  title={Presence dans les sujets analyses (2013-2026) - logique numerique et codage},
]
\begin{center}
\footnotesize\renewcommand{\arraystretch}{1.15}
\begin{tabular}{lccccccccccccc}
\toprule
\textbf{Annee} & 13 & 14 & 15 & 16 & 17 & 18 & 19 & 20 & 21 & 22 & 23 & 24 & 25 \\
\midrule
Codage binaire/hex & - & - & - & - & - & - & - & - & - & $\checkmark$ & $\checkmark$ & $\checkmark$ & - \\
Masquage / registre & - & - & - & - & - & - & - & - & - & $\checkmark$ & - & - & - \\
Logique comb. pure & - & - & - & - & - & - & - & - & - & - & - & - & - \\
\bottomrule
\end{tabular}
\end{center}
\smallskip
\footnotesize\textbf{Sous-themes frequents :}
codage binaire - resolution/precision - conversion hexadecimal - masquage ET/OU - nombre de bits pour coder une position - conversion decimal/binaire/hexa

\medskip
\textbf{Ordres de grandeur (reperes rapides).}
\begin{itemize}[leftmargin=1.3em]
  \item Reseau BT FR : tension phase-neutre typique \SI{230}{\volt}, tension phase-phase typique \SI{400}{\volt}, frequence \SI{50}{\hertz}.
  \item Courants d'etude frequents en sujet : de quelques \si{\ampere} a quelques dizaines de \si{\ampere}.
  \item En conversion d'energie, un rendement attendu reste generalement entre 0{,}8 et 0{,}98 selon la technologie.
\end{itemize}

\medskip
\textbf{Synthese express (a refaire sans support).}
\begin{enumerate}[leftmargin=1.4em]
  \item Poser les hypotheses du modele et les conventions de signe.
  \item Ecrire 2 a 4 relations cles avec unites.
  \item Realiser un calcul numerique complet et verifier l'ordre de grandeur.
  \item Justifier le resultat par un argument technique de niveau CAPET.
\end{enumerate}

\end{tcolorbox}

\CAPETENDFILE









